% GENERATED FILE. Edit canonical lessons, then run npm run generate:books.
% canonical-source-hash: fnv1a64:0f1e61c547f9256a
% canonical-lessons: BN-C267-giye, BN-C267-prothome-tarpor, BN-C267-hothat, BN-W267-letter-khanda-ta, BN-C267-ek-shonibar, BN-R267-a-story-in-order

\chapter{First, Then, Suddenly}
\label{ch:bn-first-then-suddenly}
\hlchaptermodality{267}

\begin{chapteropening}
\textbf{By the end of this chapter:} \emph{I can tell what happened in order with prôthôme, tārpôr and sheshe, join actions with the -e form (giye, kore), read a short story, and read and write \bn{ৎ} in \bn{হঠাৎ}.}

Complete BN-R267-a-story-in-order: tell a short story in order, after reading Mira's Saturday.
\end{chapteropening}

\section[\texorpdfstring{\bn{বাড়ি} \bn{গিয়ে} \bn{খেলাম} (bāṛi giye khelām)}{বাড়ি গিয়ে খেলাম (bāṛi giye khelām)}]{\bn{বাড়ি} \bn{গিয়ে} \bn{খেলাম} (bāṛi giye khelām) — one action, then the next}

\label{lesson:BN-C267-giye}



\begin{quote}
Say \textquotedblleft{}I was reading when Mira came\textquotedblright{}. (\emph{jôkhon āmi bôi pôṛchhilām, tôkhon mīrā elô}.) Then the word for \textquotedblleft{}please\textquotedblright{}, \textbf{\bn{দয়া} \bn{করে}}: \textquotedblleft{}having done a kindness\textquotedblright{}. Its \textbf{\bn{করে}} is the form of this lesson.
\end{quote}

\subsection*{You'll want to know: the -e form}
\begin{quote}
\textbf{\bn{আমি} \bn{স্টেশনে} \bn{গিয়ে} \bn{টিকিট} \bn{কিনলাম।}}
\end{quote}

\begin{quote}
\emph{āmi sṭeshône giye ṭikiṭ kinlām.} — \textquotedblleft{}I went to the station and bought a ticket.\textquotedblright{}
\end{quote}

\textbf{\bn{গিয়ে}} is \textbf{\bn{যাওয়া}} in the \emph{-e} form: \textquotedblleft{}having gone\textquotedblright{}. It finishes one action so the next can start, and the last verb carries the tense for both. Most verbs add \textbf{-\bn{এ}} to the stem: \textbf{\bn{করা}} $\to$ \textbf{\bn{করে}} (\emph{kore}), \textbf{\bn{দেখা}} $\to$ \textbf{\bn{দেখে}} (\emph{dekhe}), \textbf{\bn{বসা}} $\to$ \textbf{\bn{বসে}} (\emph{bose}). An \textbf{\bn{আ}} in the stem becomes \textbf{\bn{এ}}: \textbf{\bn{আসা}} $\to$ \textbf{\bn{এসে}} (\emph{ese}), \textbf{\bn{নামা}} $\to$ \textbf{\bn{নেমে}} (\emph{neme}), \textbf{\bn{হাঁটা}} $\to$ \textbf{\bn{হেঁটে}} (\emph{hẽṭe}). An \textbf{\bn{এ}} said \emph{e}, or an \textbf{\bn{ও}}, rises: \textbf{\bn{কেনা}} $\to$ \textbf{\bn{কিনে}} (\emph{kine}). The verbs whose stem ends in a vowel take \textbf{-\bn{য়ে}}: \textbf{\bn{খাওয়া}} $\to$ \textbf{\bn{খেয়ে}} (\emph{kheye}), \textbf{\bn{নেওয়া}} $\to$ \textbf{\bn{নিয়ে}} (\emph{niye}), \textbf{\bn{দেওয়া}} $\to$ \textbf{\bn{দিয়ে}} (\emph{diye}).

The \emph{-e} form can also say how: \textbf{\bn{হেঁটে} \bn{গেলাম}} (\emph{hẽṭe gelām}), \textquotedblleft{}I went on foot\textquotedblright{}. And it pairs with a second verb that adds a shade rather than a new action: \textbf{\bn{বসে} \bn{পড়লাম}} (\emph{bose pôṛlām}), \textquotedblleft{}I sat down\textquotedblright{}, with \textbf{\bn{পড়া}} in its other sense, \textquotedblleft{}to fall\textquotedblright{}; \textbf{\bn{হারিয়ে} \bn{গেল}}, \textquotedblleft{}it got lost\textquotedblright{}, which you know.

\subsection*{You'll want to know: one doer}
Both actions belong to the same person: \textbf{\bn{আমি} \bn{বাড়ি} \bn{গিয়ে} \bn{খেলাম}} (\emph{āmi bāṛi giye khelām}), \textquotedblleft{}I went home and ate\textquotedblright{}. When someone else does the second thing, Bengali says two sentences: \textbf{\bn{আমি} \bn{বাড়ি} \bn{গেলাম।} \bn{মীরা} \bn{খেল।}}

\subsection*{Your turn}
Say these aloud:

\begin{itemize}
\raggedright
  \item \textquotedblleft{}I got off the bus and walked home\textquotedblright{} — \emph{āmi bās theke neme hẽṭe bāṛi gelām}
  \item \textquotedblleft{}She ate and went out\textquotedblright{} — \emph{she kheye bāire gelô}
  \item \textquotedblleft{}We sat and drank tea\textquotedblright{} — \emph{āmrā bose chā khelām}
\end{itemize}

\begin{tcolorbox}[breakable,colback=teal!4,colframe=teal!35!black,title={Before you move on}]
\textbf{\bn{যাওয়া}} in the \emph{-e} form? (\textbf{\bn{গিয়ে}}.) \textbf{\bn{খাওয়া}}? (\textbf{\bn{খেয়ে}}.) Which verb carries the tense? (The last one.) Can the two actions have two doers? (No: then it is two sentences.)
\end{tcolorbox}
\section[\texorpdfstring{\bn{প্রথমে} … \bn{তারপর} … … (prôthôme … …)}{প্রথমে … তারপর … … (prôthôme … …)}]{\bn{প্রথমে} … \bn{তারপর} … \bn{শেষে} (prôthôme … tārpôr … sheshe) — the order of a story}

\label{lesson:BN-C267-prothome-tarpor}



\begin{quote}
Say \textquotedblleft{}I went home and ate\textquotedblright{}. (\emph{āmi bāṛi giye khelām}.) Then say \textbf{\bn{প্রথমে}}, \textquotedblleft{}at first\textquotedblright{}, and \textbf{\bn{শেষ}}, \textquotedblleft{}the end\textquotedblright{}.
\end{quote}

\subsection*{You'll want to know: first, then}
\begin{quote}
\textbf{\bn{প্রথমে} \bn{আমরা} \bn{চা} \bn{খেলাম।} \bn{তারপর} \bn{বাজারে} \bn{গেলাম।}}
\end{quote}

\begin{quote}
\emph{prôthôme āmrā chā khelām. tārpôr bājāre gelām.} — \textquotedblleft{}First we had tea. After that we went to the market.\textquotedblright{}
\end{quote}

\textbf{\bn{প্রথমে}}, which you know, opens a story. \textbf{\bn{তারপর}} (\emph{tārpôr}), \textquotedblleft{}after that, then\textquotedblright{}, is new: \textbf{\bn{তার}}, \textquotedblleft{}of that\textquotedblright{}, and \textbf{\bn{পর}}, \textquotedblleft{}after\textquotedblright{}, the word inside \textbf{\bn{পরে}}. \textbf{\bn{তারপর}} starts each new step; \textbf{\bn{পরে}}, \textquotedblleft{}later\textquotedblright{}, says the step came some time after.

\subsection*{You'll want to know: in the end}
\textbf{\bn{শেষে}} (\emph{sheshe}) is \textbf{\bn{শেষ}} with \textbf{-\bn{এ}}: \textquotedblleft{}at the end, finally\textquotedblright{}. It closes the story.

\begin{quote}
\textbf{\bn{শেষে} \bn{আমরা} \bn{বাড়ি} \bn{ফিরলাম।}}
\end{quote}

\begin{quote}
\emph{sheshe āmrā bāṛi phirlām.} — \textquotedblleft{}In the end we came back home.\textquotedblright{}
\end{quote}

\subsection*{Your turn}
Say these aloud:

\begin{itemize}
\raggedright
  \item \textquotedblleft{}First I went to the station by bus\textquotedblright{} — \emph{prôthôme āmi bāse sṭeshône gelām}
  \item \textquotedblleft{}After that I bought a ticket\textquotedblright{} — \emph{tārpôr ṭikiṭ kinlām}
  \item \textquotedblleft{}In the end we came back home\textquotedblright{} — \emph{sheshe āmrā bāṛi phirlām}
\end{itemize}

\begin{tcolorbox}[breakable,colback=teal!4,colframe=teal!35!black,title={Before you move on}]
What opens a story? (\textbf{\bn{প্রথমে}}.) What starts each next step? (\textbf{\bn{তারপর}}.) What closes it? (\textbf{\bn{শেষে}}.)
\end{tcolorbox}
\section[\texorpdfstring{\bn{হঠাৎ} (hôṭhāt)}{হঠাৎ (hôṭhāt)}]{\bn{হঠাৎ} (hôṭhāt) — suddenly}

\label{lesson:BN-C267-hothat}



\begin{quote}
Before the new one: say \textquotedblleft{}first\textquotedblright{}, \textquotedblleft{}after that\textquotedblright{} and \textquotedblleft{}in the end\textquotedblright{}. (\emph{prôthôme}, \emph{tārpôr}, \emph{sheshe}.)
\end{quote}

\subsection*{You'll want to know: \bn{হঠাৎ}}
\textbf{\bn{হঠাৎ}} — \emph{hôṭhāt} — \textquotedblleft{}suddenly\textquotedblright{}.

\begin{quote}
\textbf{\bn{আমি} \bn{হাঁটছিলাম}, \bn{হঠাৎ} \bn{ঝড়} \bn{এল।}}
\end{quote}

\begin{quote}
\emph{āmi hā̃ṭchhilām, hôṭhāt jhôṛ elô.} — \textquotedblleft{}I was walking, and suddenly a storm came.\textquotedblright{}
\end{quote}

A story goes \textbf{\bn{প্রথমে}}, \textbf{\bn{তারপর}}; \textbf{\bn{হঠাৎ}} is where it turns. It comes before the event that breaks into the scene.

\begin{cousinweb}
\textbf{\bn{হঠাৎ}} (\emph{hôṭhāt}) is Sanskrit \emph{haṭhāt}, \textquotedblleft{}by force, forcibly\textquotedblright{}, the ablative of \emph{haṭha}, \textquotedblleft{}force\textquotedblright{} (Monier-Williams). Monier-Williams also gives \emph{haṭhāt} as \textquotedblleft{}necessarily, inevitably\textquotedblright{}. The Bengali sense, \textquotedblleft{}suddenly\textquotedblright{}, is not among the senses it lists, and the source does not say how the word came to it.
\end{cousinweb}

\subsection*{Your turn}
Say these aloud:

\begin{itemize}
\raggedright
  \item \emph{hôṭhāt}
  \item \textquotedblleft{}I was walking, and suddenly a storm came\textquotedblright{} — \emph{āmi hā̃ṭchhilām, hôṭhāt jhôṛ elô}
  \item \textquotedblleft{}We were eating; suddenly Mira came\textquotedblright{} — \emph{āmrā khāchchhilām, hôṭhāt mīrā elô}
\end{itemize}

\begin{tcolorbox}[breakable,colback=teal!4,colframe=teal!35!black,title={Before you move on}]
What does \textbf{\bn{হঠাৎ}} mean? (\textquotedblleft{}Suddenly\textquotedblright{}.) What comes after it, a scene or an event? (An event, in the plain past.)
\end{tcolorbox}
\section[\texorpdfstring{\bn{ৎ} (t)}{ৎ (t)}]{\bn{ৎ} (t) — the letter at the end of \bn{হঠাৎ}}

\label{lesson:BN-W267-letter-khanda-ta}



\begin{quote}
Before the new one: say \textbf{\bn{হঠাৎ}}, \textquotedblleft{}suddenly\textquotedblright{}. Then read \textbf{\bn{ত}} (\emph{tô}) and the hasanta \textbf{\bn{্}}, the mark that takes a letter's vowel away.
\end{quote}

\subsection*{Script you'll notice: \bn{ৎ}}
\textbf{\bn{ৎ}} — \emph{t}, the \emph{khanda ta}, \textquotedblleft{}a piece of ta\textquotedblright{}. It is a \textbf{\bn{ত}} with no vowel after it. Where a word ends in that bare \emph{t}, as \textbf{\bn{হঠাৎ}} does, Bengali often writes \textbf{\bn{ৎ}} instead of \textbf{\bn{ত}} with a hasanta.

In the print of this book, \textbf{\bn{ৎ}} has no headline. It is a round closed loop at the top, with a long slanting tail joined to the left side of the loop; the tail's lower end, at the bottom right, is a short upright foot.

\subsection*{Writing: \bn{ৎ} — copy the printed shape}
Put your pen beside \textbf{\bn{ৎ}} and copy the shape you can see, slowly and large. Check the copy against the print: the round loop, the long tail, the short foot.

\begin{quote}
This book does not yet tell you where to start \bn{ৎ} or which way the pen goes. That is not written down here until it can be written down with a source. Copying the printed shape needs no such source, and it is how the shape becomes yours in the meantime.
\end{quote}

\subsection*{Your turn}
\begin{itemize}
\raggedright
  \item \emph{Find:} \bn{ৎ} at the end of \textbf{\bn{হঠাৎ}}
  \item \emph{Write it:} \bn{ত}, then \bn{ৎ}, from the print, naming each
  \item \emph{Write it:} \textbf{\bn{হঠাৎ}}, with \bn{ঠ} in its middle
  \item \emph{Read it:} \textbf{\bn{হঠাৎ}} aloud: \emph{hôṭhāt}
\end{itemize}

\begin{tcolorbox}[breakable,colback=teal!4,colframe=teal!35!black,title={Before you move on}]
Which letter ends \textbf{\bn{হঠাৎ}}? (\textbf{\bn{ৎ}}.) What does it say? (A \emph{t} with no vowel after it.) Does it have a headline? (No.)
\end{tcolorbox}
\section[(reading)]{(reading) — a Saturday at Grandmother's village}

\label{lesson:BN-C267-ek-shonibar}



\begin{quote}
Ninety-eight words, every one met before or explained after the story. The writer is Mira, telling of a Saturday at her grandmother's village.
\end{quote}

\begin{tcolorbox}[breakable,skin=enhanced,colback=orange!6,colframe=orange!45!black,boxrule=0.5pt,arc=1mm,left=6pt,right=6pt,top=4pt,bottom=4pt,fonttitle=\bfseries,title={Read this}]
\begin{quote}
\bn{গত} \bn{শনিবার} \bn{আমি} \bn{দিদিমার} \bn{গ্রামে} \bn{গেলাম।} \bn{প্রথমে} \bn{বাসে} \bn{স্টেশনে} \bn{গেলাম।} \bn{তারপর} \bn{ট্রেনে} \bn{চড়লাম।} \bn{ট্রেনে} \bn{আমি} \bn{একটা} \bn{বই} \bn{পড়ছিলাম।} \bn{হঠাৎ} \bn{ঝড়} \bn{এল।} \bn{ঠিক} \bn{তখন} \bn{ট্রেন} \bn{একটা} \bn{নদী} \bn{পার} \bn{হচ্ছিল।} \bn{জানালার} \bn{বাইরে} \bn{আকাশ} \bn{কালো} \bn{ছিল}, \bn{আর} \bn{গাছগুলো} \bn{কাঁপছিল।} \bn{বিকেলে} \bn{ট্রেন} \bn{গ্রামের} \bn{স্টেশনে} \bn{থামল।} \bn{তখন} \bn{ঝড়} \bn{ছিল} \bn{না।} \bn{আমি} \bn{ট্রেন} \bn{থেকে} \bn{নেমে} \bn{হেঁটে} \bn{বাড়ি} \bn{গেলাম।} \bn{বাড়ি} \bn{গিয়ে} \bn{দেখলাম}, \bn{দিদিমা} \bn{রান্না} \bn{করছিলেন।} \bn{দাদু} \bn{বারান্দায়} \bn{বসে} \bn{চা} \bn{খাচ্ছিলেন।} \bn{সন্ধ্যায়} \bn{আমরা} \bn{ভাত} \bn{আর} \bn{মাছ} \bn{খেলাম।} \bn{তারপর} \bn{দিদিমা} \bn{আমাকে} \bn{একটা} \bn{পুরনো} \bn{গল্প} \bn{বললেন।} \bn{রাতে} \bn{আবার} \bn{ঝড়} \bn{এল}, \bn{কিন্তু} \bn{আমরা} \bn{ভিতরে} \bn{ছিলাম।} \bn{শেষে} \bn{আমি} \bn{তাড়াতাড়ি} \bn{শুলাম।} \bn{কাল} \bn{সকালে} \bn{আমি} \bn{আবার} \bn{শহরে} \bn{ফিরব।} \bn{দিনটা} \bn{আমার} \bn{ভালো} \bn{লাগল।}
\end{quote}

\emph{Read it:} the passage once without stopping; do not translate

\emph{Read it:} the passage again, and sort every verb: event, scene or plan
\end{tcolorbox}

\begin{culture}
New forms: \textbf{\bn{দিদিমার}}, \textquotedblleft{}Grandmother's\textquotedblright{}, with the \textbf{-\bn{র}} of \textbf{\bn{আমার}}; \textbf{\bn{শুলাম}} (\emph{shulām}), \textbf{\bn{শোয়া}} in the past, \textquotedblleft{}I lay down to sleep\textquotedblright{}; \textbf{\bn{দিনটা} \bn{আমার} \bn{ভালো} \bn{লাগল}}, \textquotedblleft{}the day felt good to me\textquotedblright{}.

The scenes take \textbf{-\bn{ছিল}-}; the events take the plain past, joined with \emph{-e} forms and kept in order by \textbf{\bn{প্রথমে}}, \textbf{\bn{তারপর}} and \textbf{\bn{শেষে}}. \textbf{\bn{হঠাৎ}} turns the story; \textbf{\bn{ফিরব}} is the plan.
\end{culture}

\subsection*{Your turn}
Say these aloud:

\begin{itemize}
\raggedright
  \item what Mira was doing on the train when the storm came
  \item what Grandfather was doing on the veranda
\end{itemize}

\begin{tcolorbox}[breakable,colback=teal!4,colframe=teal!35!black,title={Before you move on}]
Why is \textbf{\bn{কাঁপছিল}} in the scene form? (The shaking went on through the storm.) Which word turns the story? (\textbf{\bn{হঠাৎ}}.)
\end{tcolorbox}
\section[(review)]{Review — a story in order}

\label{lesson:BN-R267-a-story-in-order}



\begin{quote}
Say \textquotedblleft{}I went to the station and bought a ticket\textquotedblright{}. (\emph{āmi sṭeshône giye ṭikiṭ kinlām}.) Then \textquotedblleft{}after that\textquotedblright{}. (\emph{tārpôr}.)
\end{quote}

\subsection*{Your turn}
\begin{itemize}
\raggedright
  \item \emph{Say it:} First I went to the village by bus.
  \item \emph{Say it:} After that I got off the bus and walked to the river.
  \item \emph{Say it:} I was sitting by the river. Suddenly Mira came.
  \item \emph{Say it:} We ate and talked. In the end I came back home.
  \item \emph{Say it:} Mira's grandmother told a story. What happened in it?
  \item \emph{Write it:} \textbf{\bn{হঠাৎ}}, then \textbf{\bn{শেষে}}
\end{itemize}

\begin{tcolorbox}[breakable,colback=teal!4,colframe=teal!35!black,title={Before you move on}]
\textbf{\bn{আসা}} in the \emph{-e} form? (\textbf{\bn{এসে}}.) What opens, steps and closes a story? (\textbf{\bn{প্রথমে}}, \textbf{\bn{তারপর}}, \textbf{\bn{শেষে}}.) Which letter ends \textbf{\bn{হঠাৎ}}? (\textbf{\bn{ৎ}}.)
\end{tcolorbox}
